\documentclass[a4paper]{article}
\usepackage[utf8]{inputenc}
\usepackage[english]{babel}
\usepackage{datetime}
\usepackage{hyperref}
\usepackage[a4paper, margin=2.5cm]{geometry}

\newdateformat{monthyeardate}{\monthname[\THEMONTH] of \THEYEAR}

\begin{document}

\thispagestyle{empty}

Francisco Figari - Buenos Aires based computer scientist

\textit{Resume} - \monthyeardate\today

\begin{itemize}

  \item Technical Manager
  \begin{itemize}
    \item Ability to simultaneously articulate technical and product needs into a coherent vision without losing the specificity of each domain.
    \item Capable of coordinating with clients to understand real needs and plan meaningful backlog which prioritizes outcomes.
    \item Ability to understand capacity, needs and blockers of development team. Ability to build collaboration between team members.
    \item Ability to steer product towards a direction which encompasses both client and internal expectations.
    \item Skilled in planning development projects and breaking them down into small, deployable steps that prioritize delivery.
    \item Capable of conducting discovery phases to assess the feasibility of new feature implementations.
    \item Experienced interdisciplinary collaborator: worked with urban planners, architects, neuroscientists, researchers, as well as frontend and backend developers.
    \item Leadership experience managing small development teams (2 to 3 developers).
    \item Familiar with agile methodologies such as Kanban and Scrum.
    \item Native Spanish speaker with advanced French (schooling) and advanced English (work).
  \end{itemize}

  \item Builder
  \begin{itemize}
    \item Ability to approach open-ended problems and create practical implementations.
    \begin{itemize}
      \item Development of investment related features like CEDEARs, argentinian stocks and time deposits.
      \item Research and development done in the context of neuro-science regarding the feasibility of web-based \texttt{eye-tracking} applied to remote clinical trials.
      \item Knowledgeable in urban planning topics. Development of a procedural generator for affordable multi-family housing adhering to strict urban planning codes. Development of real estate and housing simulation.
    \end{itemize}
    \item Capable of working within medium-scale microservice (\texttt{ms}) ecosystems.
    \begin{itemize}
      \item Maintenance and integration of new \texttt{ms} into the network, both consumers and producers.
      \item Implementation of new features spanning multiple \texttt{ms}.
      \item General knowledge of tracing and monitoring tools (e.g., \texttt{Datadog}).
      \item Scope: 10 services, with the largest handling 200 requests per second during market hours.
    \end{itemize}
    \item Skilled in translating problems described in natural language into formal representations.
    \item Ability to scale codebases, perform refactoring, and decompose complex function flows into smaller, manageable subflows.
    \item General understanding of multiple programming paradigms: object-oriented, procedural, functional, and event-oriented.
    \item Professional experience with multiple programming languages: \texttt{Typescript}, \texttt{Golang}, \texttt{Python}, \texttt{C++}.
    \item Designs with computational complexity in mind (Big-\texttt{O} notation).
    \item Capable of presenting ideas clearly and concisely, following formats similar to scientific papers.
    \item Skilled in defining interfaces that involve multiple technologies (e.g., backend-to-mobile communication).
    \item Capable of designing software flows based on HTTP resources.
    \item General knowledge of data persistence technologies (\texttt{PostgreSQL}, \texttt{Redis}).
    \item Conscious and effective use of \texttt{git} in development workflows.
    \item General understanding of \texttt{Unix} systems.
    \item Effective usage of AI agents (e.g. \texttt{claude-code}) to accelerate development processes.
  \end{itemize}

  \item Mentor
  \begin{itemize}
    \item Capable of reviewing coworkers work with attention to quality of the proposed solution.
    \item Capable of tracking junior developers’ progress over time in order to guide them and teach better development practices.
    \item Ability to teach numerical methods. Experienced in evaluating practical assignments with emphasis on clarity of idea presentation.
  \end{itemize}

% Separator line
\vspace{1em}
\noindent\rule{\linewidth}{0.4pt}
\vspace{1em}

% Work History section
\item Work History
\begin{itemize}
  \item Senior software engineer @ \textit{Urbanly} since August 2025.
  \item L4 backend engineer @ \textit{Brubank} from January 2025 to August 2025.
  \item L3 backend engineer @ \textit{Brubank} from Apr 2023 to January 2025.
  \item Research intern @ \textit{DC, FCEyN, UBA} between August 2021 and September 2022 under scholarship program for
    initiation in research. Work done around thesis project.
  \item Fullstack programmer @ \textit{DubbingDigital} between March 2021 y December 2021.
  \item Lead programmer @ \textit{UrbanSim} between July 2018 y November 2020.
  \item Research intern @ \textit{ARSLab, Carleton University} between February and June of 2018.
  \item Teaching assistant in the Numerical Methods course @ \textit{DC, FCEyN, UBA}, 4 semesters between 2016 y 2018.
  \item Freelance programmer @ \textit{LIAA, FCEyN, UBA} at the end of 2017.
  \item Junior programmer @ \textit{Eryx} during 4 months at the beginning of 2016.
\end{itemize}

% Education section
\item Education and Training
\begin{itemize}
    \item Product Discovery for Agile Teams @ \textit{Sense and Respond Learning}; 8 hours, intro to \textit{Lean Product Canvas}
    \item Yoga teaching training program @ \textit{Red Prisma} between July 2023 and April 2025. 18 intensive weekends plus weekly theoretical classes.
    \item \textit{Licenciatura} in Computer Science @ \textit{DC, FCEyN, UBA}, equivalent to Bachelor + Master with thesis, started in 2013 and finished in 2022.
    \item High school @ \textit{Licée Jean Mermoz, Buenos Aires}, finished in 2013.
\end{itemize}

\end{itemize}
\end{document}
